\section{Fences, socles, and coordinate transport}
\label{sec:framework}

All groups in this section are finite. A vertex of a subgroup interval is
written as a capital letter and identified with the corresponding subgroup.
The common bottom is denoted by $D$. We first record the normal-subgroup
test that will give the vertices their simple-group labels. We then explain
the two different kinds of factor transport: diagonal merging when the
simple type is unchanged, and coordinate products when it changes.

\subsection{The normal-subgroup test}

\begin{definition}
A bounded lattice interval $[a,b]$ is \emph{fenced} if it has an interior
element and, for each $a<c<b$, there are $a<u<v<b$ such that
\[
 c\wedge u=c\wedge v=a,
 \qquad c\vee u=c\vee v=b.
\]
Thus the two comparable elements $u,v$ are both complements of $c$ in
the interval. The requirement of an interior element excludes vacuous
fences on covering intervals. Length always counts edges.
\end{definition}

\begin{lemma}[Normal-subgroup test]
\label{frm:normal-test}
Suppose that $[D,X]$ is fenced. If $R\trianglelefteq X$, then either
$R\le D$ or $DR=X$. In particular, $[D,X]$ is not modular.
The same assertion applies to a normal subgroup of $X/F$ whenever
$F\trianglelefteq X$ and $F\le D$, with bottom $D/F$.
\end{lemma}

\begin{proof}
For $D\le U\le V\le X$, the normality of $R$ and the subgroup
modular identity give
\begin{equation}
 \label{frm:normal-modular-identity}
 (DR\vee U)\wedge V
   =(RU)\cap V
   =U(R\cap V)
   =U\vee(DR\wedge V).
\end{equation}
If $D<DR<X$, apply the fence to $DR$ and its comparable complements
$U<V$. Equation~\eqref{frm:normal-modular-identity} would give $V=U$.
Thus $DR$ is an endpoint. In a modular lattice the same identity holds
with any element in place of $DR$, so the existence of a fence also
excludes modularity. Finally, taking inverse images identifies
$[D/F,X/F]$ with $[D,X]$.
\end{proof}

We will repeatedly use the following explicit form of the subgroup
correspondence. If $R\trianglelefteq X$ and $X=DR$, then
\begin{equation}
 \label{frm:supplement-correspondence}
 [D,X]\longrightarrow
 \{E\le R:D\cap R\le E,\ E\text{ is normalized by }D\},
 \qquad U\longmapsto U\cap R
\end{equation}
is an order isomorphism, with inverse $E\mapsto DE$. Indeed
$U=D(U\cap R)$ for $D\le U\le DR$, and $(DE)\cap R=E$ when
$D\cap R\le E$. This correspondence includes every overgroup of $D$.

\begin{lemma}[Labels]
\label{frm:labels}
Suppose that $[D,X]$ is fenced, and put
\[
 F_X=\operatorname{core}_X(D)
     =\bigcap_{g\in X}D^g.
\]
Then $X/F_X$ has a unique minimal normal subgroup
\begin{equation}
 \label{frm:socle-label}
 S_X=N_X/F_X\cong T_X^{m_X},
\end{equation}
where $T_X$ is nonabelian simple and $N_X$ is its full inverse image in
$X$. Moreover,
\[
 X=DN_X.
\]
Every nontrivial normal subgroup of $X/F_X$ contains $S_X$, and
\[
 C_{X/F_X}(S_X)=1.
\]
In particular, $X/F_X$ has no nontrivial soluble normal subgroup.
If $I_X$ is the set of simple direct factors of $S_X$, then $D$ and $X$
induce the same transitive permutation group on $I_X$.
\end{lemma}

\begin{proof}
Factor out $F_X$, so that the action of $X$ on the cosets of $D$ is
faithful. Lemma~\ref{frm:normal-test} says that every nontrivial normal
subgroup supplements $D$, or equivalently is transitive in this action.
Let $M$ be a minimal normal subgroup. If $M$ were abelian,
\eqref{frm:supplement-correspondence} would identify $[D,X]$ with an
interval of $D$-invariant subgroups of an abelian group. The latter lattice
is modular: intersection is its meet and subgroup product is its join.
This contradicts Lemma~\ref{frm:normal-test}.

If $M'$ were another minimal normal subgroup, then
$[M,M']\le M\cap M'=1$. Every $D$-invariant subgroup of $M$ would
therefore be normalized by $DM'=X$, and so would be either $1$ or $M$.
The correspondence~\eqref{frm:supplement-correspondence} would leave
at most two elements in $[D,X]$, again a contradiction.

A finite minimal normal subgroup is characteristically simple, hence
a direct product of isomorphic simple groups. Since $M$ is nonabelian,
these simple groups are nonabelian. Conjugation by $X$ is transitive on
the direct factors: otherwise the product of one orbit of factors would
be a proper nontrivial normal subgroup of $X$ contained in $M$. This
proves~\eqref{frm:socle-label}. Every nontrivial normal subgroup contains
a minimal normal subgroup and therefore contains $M$. Its centralizer
is normal and intersects $M$ in $Z(M)=1$, so that centralizer is trivial.
The assertion about soluble normal subgroups follows as well.

Finally $X=DM$. Inner automorphisms of $M$ fix each of its simple direct
factors. Thus $D$ and $X$ have the same image on their factor set.
Restoring $F_X$ proves the statements as written.
\end{proof}

We call an $X$ satisfying Lemma~\ref{frm:labels} a \emph{labelled
vertex}, and call the isomorphism type of $T_X$ its \emph{simple type}.
The preceding argument is also the strongly nonmodular normal-subgroup
argument in \cite[Lemma~3.6]{PalfyTwisted}.

\begin{corollary}[Bottlenecks]
\label{frm:bottleneck}
Let $D<B\le Z$, suppose that $[D,B]$ is fenced, and suppose that $B$
is comparable with every member of $[D,Z]$. If $R\trianglelefteq Z$
and $R\nleq D$, then
\[
 B=D(R\cap B).
\]
Consequently a soluble normal subgroup of $Z$ is contained in $D$.
If $R,Q\trianglelefteq Z$ and
\[
 [R,Q]\le F\trianglelefteq Z,\qquad F\le D,
\]
then at least one of $R,Q$ is contained in $D$.
These conclusions also hold with $B=Z$.
\end{corollary}

\begin{proof}
The subgroup $R\cap B$ is normal in $B$. Suppose first that
$R\cap B\le D$. If $B\le DR$, then
\[
 B=B\cap DR=D(B\cap R)=D,
\]
a contradiction. Comparability therefore gives $DR\le B$, whence
$R=R\cap B\le D$. Its contrapositive, followed by
Lemma~\ref{frm:normal-test} in $B$, gives the asserted supplementation.
The image of $R\cap B$ modulo $F_B$ then contains the nonabelian
socle $S_B$, by Lemma~\ref{frm:labels}. This is impossible if $R$ is
soluble. For the final assertion, $F\le F_B$ because $F$ is normal in
$B$ and contained in $D$. If both $R$ and $Q$ escaped $D$, their
intersections with $B$ would both map over $S_B$, while their images
would commute. This would make $S_B$ abelian.
\end{proof}

\subsection{Subdirect powers and unchanged simple types}

For a direct product $T^I=\prod_{i\in I}T_i$, with each $T_i\cong T$,
a subgroup is \emph{subdirect} if all coordinate projections are
surjective. A \emph{full diagonal strip} on a nonempty subset $J\subseteq I$
is a subgroup obtained by identifying all the $T_j$, $j\in J$, by
isomorphisms, and putting identity elements outside $J$.
We recall the standard diagonal-strip description
\cite[Lemma~2.1]{PalfyTwisted}, including its behavior under inclusion.

\begin{lemma}[Subdirect powers]
\label{frm:subdirect}
Let $T$ be nonabelian simple. Every subdirect subgroup $E\le T^I$
is a direct product of full diagonal strips on the parts of a partition
of $I$. Both the partition and the coordinate identifications within
each part are determined by $E$. An inclusion of subdirect subgroups
corresponds to a refinement of these parts, with the existing
identifications restricted to the smaller parts.
\end{lemma}

\begin{proof}
Induct on $|I|$, separating the last coordinate. The projection to the
other coordinates is, by induction, a product $T^r$ of diagonal strips.
We have a subdirect subgroup of $T^r\times T$. The kernels of its two
projections are normal in the respective factors. The kernel on the last
$T$ is either $T$, in which case this is the full product, or $1$, in
which case the subgroup identifies a quotient of $T^r$ with $T$.
Normal subgroups of $T^r$ are products of its simple direct factors, so
that quotient is one of the factors. The last coordinate is therefore
joined to exactly one existing strip.

Two coordinates belong to the same strip precisely when the projection
of $E$ to their product is the graph of an isomorphism rather than the
full product. This characterizes the partition and its links, and also
proves the assertion about inclusion.
\end{proof}

\begin{proposition}[Equal-type transport]
\label{frm:equal-transport}
Let $D<X<Y$ be labelled vertices. Then $F_Y\le F_X$, and $T_X$ is
isomorphic to a section of $T_Y$. Put
\[
 E_X=(X\cap N_Y)/F_Y\le S_Y.
\]
The following conditions are equivalent:
\begin{enumerate}[label=\textup{(\roman*)}]
 \item $T_X\cong T_Y$;
 \item $E_X$ is subdirect in $S_Y$.
\end{enumerate}
When they hold, the quotient map induces an isomorphism
\begin{equation}
 \label{frm:equal-socle-identification}
 E_X\xrightarrow{\ \sim\ }S_X,
 \qquad
 N_X=F_X(X\cap N_Y),
 \qquad F_X\cap N_Y=F_Y.
\end{equation}
Moreover $m_X<m_Y$, and the strip supports give a canonical
$X$-equivariant surjection
\begin{equation}
 \label{frm:equal-support-map}
 \sigma_{YX}:I_Y\longrightarrow I_X.
\end{equation}
For three comparable labelled vertices of the same simple type these
maps compose:
$\sigma_{ZX}=\sigma_{YX}\sigma_{ZY}$.
\end{proposition}

\begin{proof}
The normal subgroup $F_Y$ of $X$ is contained in $D$, so $F_Y\le F_X$.
Since $Y=DN_Y$ and $D\le X$,
\[
 X=D(X\cap N_Y).
\]
The image of $E_X$ in $X/F_X$ is a nontrivial normal subgroup and
therefore contains $S_X$. Every simple composition factor of a subgroup
of a direct product is a section of one of the factors: projection to
one coordinate gives an exact sequence whose kernel lies in the product
of the remaining coordinates, and induction applies. Thus $T_X$ is a
section of $T_Y$.

If $T_X\cong T_Y$, some coordinate projection of $E_X$ is all of $T_Y$;
otherwise none of those smaller groups could have a composition factor
of order $|T_Y|$. The transitive action of $D$ on $I_Y$ makes all the
projections surjective. Conversely, suppose $E_X$ is subdirect.
Lemma~\ref{frm:subdirect} expresses it as a product of diagonal copies
of $T_Y$. The action of $D$ on those copies is transitive. Since
$X/F_Y=(D/F_Y)E_X$, the subgroup $E_X$ is minimal normal in $X/F_Y$.
Its intersection with $F_X/F_Y$ is either trivial or all of $E_X$.
The latter would imply $X=D$, which is excluded. Its image is therefore
the unique minimal normal subgroup $S_X$, proving the equivalence and
\eqref{frm:equal-socle-identification}.

There are $m_X$ strips on $m_Y$ coordinates. If $m_X=m_Y$, every strip
is a singleton, so $E_X=S_Y$ and $X=Y$. Thus $m_X<m_Y$.
The quotient identification with $S_X$ assigns to every strip its
actual simple factor in $I_X$ and defines~\eqref{frm:equal-support-map}.
The subgroup $E_X$ is normal in $X/F_Y$, and its quotient identification
with $S_X$ commutes with conjugation by $X$; hence this support map is
$X$-equivariant.

For composition, factor out $F_Z$ and let $E_Y=Y\cap N_Z$.
We have $N_Y=F_YE_Y$ and $F_Y\cap E_Y=1$. Since $F_Y\le X$,
\[
 X\cap N_Y=F_Y(X\cap E_Y)=F_Y(X\cap N_Z).
\]
Thus the subdirect subgroup used for $X<Y$ is the image of the one
used for $X<Z$ under the already specified identification
$E_Y\cong S_Y$. Refinement of its strip partition is exactly
composition of the support maps.
\end{proof}

In particular, a change of simple type is strict in order:
$T_X\not\cong T_Y$ implies $|T_X|<|T_Y|$. A section of a finite group
having its full order must be the group itself.

\subsection{Products and one-coordinate intervals}

Fix a labelled upper vertex $Y$. Work modulo $F_Y$ when discussing its
socle $S_Y=T^{I_Y}$. Conjugation by $D$ is transitive on $I_Y$.
Choose a reference coordinate and transport identifications of the other
coordinates by elements of $D$. Notation such as $P^{I_Y}$ always means
the product of these transported copies. The subgroup in the reference
coordinate must be invariant under the corresponding factor stabilizer;
this is precisely the condition that makes the transported product
invariant under the whole group in question.

Call $D\le Z\le Y$ \emph{saturated relative to $Y$} if
$(Z\cap N_Y)/F_Y$ is the product of its coordinate projections.
The proof below uses the product--subdirect intersection arguments of
\cite[Lemmas~2.12 and~2.13]{PalfyTwisted} to obtain the required saturation.

\begin{proposition}[Strict products and saturation]
\label{frm:strict-product}
Let $D\le X<Y$, where $Y$ is labelled. Suppose that the coordinate
projections of
\[
 B=(X\cap N_Y)/F_Y
\]
are proper and nontrivial. If $[X,Y]$ is fenced, or if $X$ is a coatom
of $[D,Y]$, then $X$ is saturated relative to $Y$.

Once $B=P^{I_Y}$ is such a product, no proper subdirect subgroup of
$S_Y$ contains $B$. Every coatom of $[X,Y]$, and every intersection
of these coatoms, is saturated relative to $Y$.
These assertions do not require $X$ to be labelled. For labelled
$X<Y$ of different simple types the projection hypothesis holds.
\end{proposition}

\begin{proof}
Factor out $F_Y$, and let $C$ be the product of the coordinate
projections of $B$. It is $X$-invariant. The correspondence
\eqref{frm:supplement-correspondence}, with bottom $X$ and normal
subgroup $S_Y$, identifies $[X,Y]$ with the $X$-invariant subgroups of
$S_Y$ containing $B$.

Suppose $B<C<S_Y$. If $U<V$ are comparable complements of $XC$,
put $E_U=U\cap S_Y$ and $E_V=V\cap S_Y$. Their projections already
contain the projections of $C$. Hence projecting
$\langle E_U,C\rangle=S_Y$ shows that $E_U$ is subdirect; the same
holds for $E_V$.

Every link within a strip of $E_U$ carries the projection of $B$ in one
coordinate onto its projection in the other, because $B\le E_U$.
Consequently $E_U\cap C$ is a product of independent diagonal copies
of those nontrivial projections, one for each strip of $E_U$.
The proper inclusion $E_U<E_V$ splits at least one strip, by
Lemma~\ref{frm:subdirect}. Intersecting with $C$ then permits an
additional independent nonidentity coordinate value, and gives
$E_U\cap C<E_V\cap C$. This contradicts the required equality of
both intersections with $B$. A fence therefore forces $B=C$.
If $X$ is a coatom, $B<C<S_Y$ directly contradicts maximality.

A proper subdirect subgroup has a link between two distinct coordinates.
The full product $P^{I_Y}$ contains an element which is identity in the
first of these coordinates and nonidentity in the second, contradicting
that link. Now let $Z$ be a coatom above $X$. If its projections were
full, it would be such a proper subdirect subgroup. Otherwise its product
closure is proper, and maximality forces it to equal that closure.
Intersections of products are products.

Finally, for different labelled simple types, the projections are
nontrivial because their image contains $S_X$, and transitivity makes
all of them nontrivial. If one were full, all would be full, contradicting
Proposition~\ref{frm:equal-transport}.
\end{proof}

\begin{proposition}[One-coordinate projection]
\label{frm:projection}
Let $D\le X<Y$, with $Y$ labelled, and suppose that
\[
 (X\cap N_Y)/F_Y=P_X^{I_Y},\qquad 1<P_X<T_Y.
\]
Suppose also that $X$ is an intersection of coatoms of $[X,Y]$;
a coatom itself is allowed. Fix $i\in I_Y$ and let
\[
 \rho:Y_i\longrightarrow\operatorname{Aut}(T_Y)
\]
be the action of its factor stabilizer on that coordinate. Identify
$T_Y$ with its inner automorphism group, and put
$H_* = \rho(X_i)$. Then
\begin{equation}
 \label{frm:inner-bottom}
 H_*\cap T_Y=P_X.
\end{equation}
The saturated overgroups of $X$ form a sublattice of $[X,Y]$, and
\begin{equation}
 \label{frm:coordinate-interval}
 Z\longmapsto K_Z=\rho(Z_i)=H_*P_Z,
 \qquad (Z\cap N_Y)/F_Y=P_Z^{I_Y},
\end{equation}
is an isomorphism from that sublattice onto the entire ordinary subgroup
interval $[H_*,H_*T_Y]$. Original covering relations between saturated
vertices remain covering relations after projection. In particular, a
strict change of simple type at a cover yields a maximal subgroup of an
almost simple group, supplementing its simple socle.
\end{proposition}

\begin{proof}
For a coatom $M$ above $X$, Proposition~\ref{frm:strict-product} writes
its intersection with $S_Y$ as $P_M^{I_Y}$, with $1<P_M<T_Y$.
Let $J=H_*\cap T_Y$. It normalizes $P_M$, so
\[
 P_MJ\le N_{T_Y}(P_M)<T_Y.
\]
The last inequality follows from simplicity and $1<P_M<T_Y$.
Both $P_M$ and $J$ are $H_*$-invariant, and their product therefore
transports to a saturated overgroup of $M$ which is still proper in $Y$.
Maximality gives $J\le P_M$. Intersecting over the coatom representation
of $X$ gives $J\le P_X$; the reverse inclusion is immediate. This proves
\eqref{frm:inner-bottom}.

Products above $P_X^{I_Y}$ correspond precisely to $H_*$-invariant
subgroups $P_Z\le T_Y$ containing $P_X$. By
\eqref{frm:inner-bottom}, these in turn correspond precisely to the
subgroups $K_Z$ of $H_*T_Y$ containing $H_*$: the two maps are
$P_Z\mapsto H_*P_Z$ and $K_Z\mapsto K_Z\cap T_Y$.
Since $Z=X P_Z^{I_Y}$ and the product fixes the coordinate set,
$Z_i=X_iP_Z^{I_Y}$, proving~\eqref{frm:coordinate-interval}.
Intersections and generated subgroups of coordinate products are again
coordinate products, proving the sublattice assertion. Any intermediate
projected subgroup would lift to an original intermediate saturated
subgroup, which proves the assertion about covers. At a strict cover
Proposition~\ref{frm:strict-product} supplies the product hypothesis,
and the bottom is itself a coatom. The ambient group $H_*T_Y$ is a
subgroup of $\operatorname{Aut}(T_Y)$ containing $T_Y$, hence is almost
simple.
\end{proof}

\subsection{The active chief factor and its core}

We will need the action on a specified nonabelian chief factor in a
projected group. It is not necessary to assert that every projected
coset action is quasiprimitive. The following proposition keeps the
original cores and the projected cores distinct.

\begin{proposition}[Active-factor transport]
\label{frm:active-transport}
Let $D<Z\le Y$ be labelled vertices, and suppose that $Z$ is saturated
relative to $Y$:
\[
 (Z\cap N_Y)/F_Y=P_Z^{I_Y}.
\]
Let $K_Z$ be the automorphism group induced by a factor stabilizer in
$Z$ on a reference coordinate of $S_Y$. There are $K_Z$-normal
subgroups $Q_Z\le R_Z\le P_Z$ such that
\begin{align}
 \label{frm:product-core}
 (F_Z\cap N_Y)/F_Y&=Q_Z^{I_Y},\\
 \label{frm:product-active-preimage}
 (N_Z\cap N_Y)/F_Y&=R_Z^{I_Y},\\
 \label{frm:coordinate-chief}
 R_Z/Q_Z&\cong T_Z^{a_Z},\qquad m_Z=a_Zm_Y.
\end{align}
Here $a_Z$ is a positive integer. The section $R_Z/Q_Z$ is a chief factor for the action of $K_Z$ on
$P_Z$: there is no $K_Z$-normal subgroup strictly between $Q_Z$ and
$R_Z$. Its factors give a $Z$-equivariant surjection
\begin{equation}
 \label{frm:product-support}
 \tau_{ZY}:I_Z\longrightarrow I_Y
\end{equation}
whose fibers have size $a_Z$. The stabilizers in $D$ and $Z$ of a
coordinate in $I_Y$ induce the same transitive action on its fiber.

In addition, suppose $X<Z$ and the hypotheses of
Proposition~\ref{frm:projection} hold for $X<Y$. Set
\[
 H_* = \rho(X_i),\qquad K_Z=H_*P_Z,
 \qquad C_Z=\operatorname{core}_{K_Z}(H_*).
\]
Then the subgroups in~\eqref{frm:product-core} satisfy
\begin{equation}
 \label{frm:projected-core}
 Q_Z=C_Z\cap P_Z,
 \qquad K_Z=H_*R_Z.
\end{equation}
Thus $R_ZC_Z/C_Z\cong R_Z/Q_Z$ is a transitive minimal normal
subgroup in the coset action of $K_Z/C_Z$ on $H_*/C_Z$.
The groups $H_*$ and $K_Z$ induce the same action on its simple factors.
No assertion that this is the entire projected socle is required.

If $a_Z=1$, the quotient $P_Z/R_Z$ is soluble.
\end{proposition}

\begin{proof}
Factor out $F_Y$ throughout the proof. Put
\[
 E=Z\cap N_Y=P_Z^{I_Y},\qquad F=F_Z,
 \qquad W=F\cap E.
\]
Let $Q_i$ be the projection of $W$ to the $i$th coordinate, and let
$\widehat Q=\prod_iQ_i$. These projections are transported into one
another by $D$, and each is invariant under its coordinate stabilizer.
Thus $\widehat Q$ is normal in $Z$.

We claim that $W=\widehat Q$. For $q,q'\in Q_i$, choose $w\in W$
whose $i$th coordinate is $q$ and let $e_i(q')\in E$ be supported only
on coordinate $i$. Since $W\trianglelefteq E$,
\[
 [w,e_i(q')]=e_i([q,q'])\in W.
\]
It follows that $[\widehat Q,\widehat Q]\le W$.
The image of $\widehat Q$ in $Z/F$ is therefore the abelian normal
subgroup $\widehat Q/W$. Lemma~\ref{frm:labels} forces that image to
be trivial. Hence $W=\widehat Q$, proving~\eqref{frm:product-core}.
This argument also covers $Z=X$ at a covering comparison; it does not
use a projected core with bottom equal to the whole projected group.

The normal subgroup $EF/F$ of $Z/F$ is nontrivial, since $Z=DE$ and
$Z>D$. It contains $S_Z$. Therefore, if $L=N_Z\cap E$, then
\[
 N_Z=FL,\qquad L/W\cong S_Z.
\]
In particular $L/W$ is a perfect normal subgroup of
$E/W=\prod_i(P_i/Q_i)$.

Every perfect normal subgroup $V$ of a direct product $\prod_iB_i$ is
the product of its intersections with the factors. To see this, for
$v,w\in V$ the element $[v,e_i(w_i)]$ lies in $V\cap B_i$ and is
the $i$th coordinate of $[v,w]$. Hence
$V'\le\prod_i(V\cap B_i)\le V$, and $V=V'$ gives equality.
Apply this observation to $L/W$. It splits as a product of its
coordinate intersections, each a direct product of copies of $T_Z$.
Taking full inverse images gives $L=R_Z^{I_Y}$ and
\eqref{frm:product-active-preimage}--\eqref{frm:coordinate-chief}.
Conjugation by $Z$ permutes these coordinate intersections and commutes
with $L/W\cong S_Z$, so the resulting support map is $Z$-equivariant.
Transitivity of $D$ on $I_Y$ makes the multiplicity $a_Z$ the same in
every coordinate.

Any $K_Z$-normal subgroup strictly between $Q_Z$ and $R_Z$ would
transport to a $Z$-normal subgroup strictly between $W$ and $L$.
Its image would be a proper nontrivial normal subgroup of $Z/F$
inside $S_Z$. This proves the chief-factor assertion and also the
transitivity of $K_Z$ on the $a_Z$ simple factors of $R_Z/Q_Z$.

We have $Z=DN_Z=DL$. Since $L\le N_Y$ fixes $I_Y$, for a chosen
coordinate $i$ this gives
\[
 Z_i=D_iL,\qquad K_Z=\rho(D_i)R_Z.
\]
The group $R_Z$ acts by inner automorphisms on $R_Z/Q_Z$ and fixes
its simple factors. Thus $\rho(D_i)$ and $K_Z$ have the same image
on that factor set. The kernel of $\rho$ centralizes the ambient
coordinate and so acts trivially on this section. This identifies the
action with the original fiber action and proves
\eqref{frm:product-support} and its assertions about stabilizers.

For the additional projected-core assertion, let
$C_Z=\operatorname{core}_{K_Z}(H_*)$. The subgroup $Q_Z$ is normal
in $K_Z$ and lies in $P_X$, because $F_Z\le D\le X$.
By~\eqref{frm:inner-bottom} it lies in $H_*$, and therefore
$Q_Z\le C_Z\cap P_Z$. Conversely, transport $C_Z\cap P_Z$ over all
coordinates. Its product is normal in $Z$ and is contained in
$P_X^{I_Y}\le X<Z$. Lemma~\ref{frm:normal-test} for the original
interval $[D,Z]$ forces this product into $D$, hence into $F_Z$.
Equation~\eqref{frm:product-core} now gives the reverse inclusion.

Since $\rho(D_i)\le H_*\le K_Z$ and
$K_Z=\rho(D_i)R_Z$, we obtain $K_Z=H_*R_Z$. Also
$R_Z\cap C_Z=Q_Z$. The resulting image of $R_Z$ modulo $C_Z$ is
minimal normal by the chief-factor assertion and is transitive by
supplementation. Its factor action is the same for $H_*$ and $K_Z$,
because it is already induced by $\rho(D_i)$.

Finally suppose $a_Z=1$. The support map~\eqref{frm:product-support}
is then a bijection. The subgroup $E$ fixes $I_Y$ and consequently fixes
$I_Z$. Conjugation embeds $EF/F$ faithfully in
$\operatorname{Aut}(S_Z)$, since $C_{Z/F}(S_Z)=1$ by
Lemma~\ref{frm:labels}. Its image therefore lies in
$\operatorname{Aut}(T_Z)^{I_Y}$ and contains the inner subgroup
$S_Z$. The quotient $E/L\cong(P_Z/R_Z)^{I_Y}$ embeds in
$\operatorname{Out}(T_Z)^{I_Y}$. Schreier's Theorem, a consequence
of the classification of finite simple groups, says that
$\operatorname{Out}(T_Z)$ is soluble; see also the explicit use in
\cite[Section~1]{PalfyTwisted}. Hence $P_Z/R_Z$ is soluble.
\end{proof}

\begin{remark}
The projected bottom in Proposition~\ref{frm:projection} is the image
of $X_i$, not necessarily the image of $D_i$ alone: in general
$H_* = \rho(D_i)P_X$. The smaller image already induces the full
active factor actions, so the actual projected bottom does also.
Likewise, the identity $Q_Z=C_Z\cap P_Z$ in
Proposition~\ref{frm:active-transport} is asserted only for $X<Z$.
The product-core and chief-factor statements preceding it apply also
to the lower endpoint of a covering comparison. These distinctions
allow us to use maximal-subgroup information without assuming more
quasiprimitivity in the projected diagram than has been proved.
\end{remark}
